% ============================================================================
% Курсов проект по дисциплина „UML обектно-ориентирано проектиране“
% ОКС „Магистър“, ФКСТ, ТУ-София
%
% Тема: инструмент за двупосочно инженерство между изходен код на C++,
%       UML диаграми на класове и релационни схеми на бази от данни.
%
% Компилация:  latexmk -pdf Main.tex        (pdfLaTeX, НЕ XeLaTeX)
% Форматиране: preamble.tex
%
% Този файл е СКЕЛЕТЪТ: структурата и редът на разделите. Съдържанието живее
% в chapters/. Незапълнените места са маркирани с \TODO{...} и се намират с
%     grep -rn 'TODO' chapters/ Main.tex references.bib
% ============================================================================
\documentclass[12pt]{article}

% ============================================================================
% preamble.tex, форматиране по изискванията на ФКСТ.
%
% Всяка секция тук съответства на правило от REQUIREMENTS.md (Част А) и на
% реда му в таблицата на превода (Част Б). Не пипай стойностите без да
% провериш срещу REQUIREMENTS.md. Те са нормативни, не са вкус.
%
% Двигател: pdfLaTeX (T2A + tempora). НЕ XeLaTeX, вж. REQUIREMENTS.md Б.0.
% ============================================================================

\usepackage{cmap}
\usepackage[T2A]{fontenc}
\usepackage[utf8]{inputenc}
\usepackage[bulgarian]{babel}

% --- Типография (А.4) ---------------------------------------------------
% tempora = Times-метрично съвместим шрифт с пълна кирилица под T2A.
\usepackage{tempora}
\usepackage[top=2.5cm, bottom=2.5cm, left=3.0cm, right=2.0cm, a4paper]{geometry}
\usepackage{setspace}
\setstretch{1.5}                     % „междуредие 1.5 lines“, калибруем, вж. Б.2
% „без пренасяния на срички“ (А.4), чрез ЗАБРАНИТЕЛНИ НАКАЗАНИЯ, не чрез
% \usepackage[none]{hyphenat}. Разликата е съществена: `hyphenat[none]` задава
% \hyphenchar=-1 и така изключва самия механизъм за прекъсване, от който зависи
% `breaklines` на listings, дългите редове в листингите тогава НЕ се пренасят и
% излизат извън полето (1227 препълвания дойдоха оттам). Наказанията постигат същия
% резултат за текста (TeX никога не пренася), но оставят механизма достъпен.
\usepackage[expansion=false]{microtype}
\DisableLigatures{encoding=*,family=tt*}
\hyphenpenalty=10000
\exhyphenpenalty=10000
\emergencystretch=1.5em              % цената на justified-без-преноси, вж. Б.4
\sloppy
\usepackage{indentfirst}             % отстъп и на ПЪРВИЯ абзац след заглавие, вж. Б.3
\setlength{\parindent}{1.25cm}
\setlength{\parskip}{0pt}

% --- Номерация на страниците: долу вдясно (А.4) -------------------------
\usepackage{fancyhdr}
\pagestyle{fancy}
\fancyhf{}
\renewcommand{\headrulewidth}{0pt}
\fancyfoot[R]{\thepage}

% --- Бележки под линия: 10 pt, single, justified (А.4) -------------------
\usepackage[hang,flushmargin]{footmisc}
\usepackage{ragged2e}
\renewcommand{\footnotelayout}{\fontsize{10}{12}\selectfont\setstretch{1.0}\justifying}

% --- Заглавия на раздели / параграфи / подпараграфи (А.7) ---------------
\usepackage{titlesec}
\newcommand{\sectionbreak}{\clearpage}   % „всеки раздел започва на нова страница“

\renewcommand{\thesection}{\Roman{section}}
\renewcommand{\thesubsection}{\arabic{section}.\arabic{subsection}}
\renewcommand{\thesubsubsection}{\arabic{section}.\arabic{subsection}.\arabic{subsubsection}}

% Раздел: римски, ALL CAPS, центриран, bold 14 pt
\titleformat{\section}[block]
  {\fontsize{14}{16.8}\bfseries\centering}{\thesection. }{0pt}{\MakeUppercase}

% Параграф: арабски, bold 14 pt, отстъп 1.25 cm
\titleformat{\subsection}[block]
  {\fontsize{14}{16.8}\bfseries}{\thesubsection. }{0pt}{}
\titlespacing*{\subsection}{1.25cm}{*3}{*2}

% Подпараграф: bold 13 pt, отстъп 2.0 cm
\titleformat{\subsubsection}[block]
  {\fontsize{13}{15.6}\bfseries}{\thesubsubsection. }{0pt}{}
\titlespacing*{\subsubsection}{2.0cm}{*2}{*1}

% --- Съдържание: разделите с ГЛАВНИ букви, 12 pt, bold (А.6) -------------
\usepackage{etoolbox}
\usepackage[titles]{tocloft}
% NB: кръпката ТРЯБВА да е СЛЕД tocloft, защото пакетът предефинира \l@section и
% иначе я изтрива (тихо: няма грешка, просто разделите излизат с малки букви).
\makeatletter
\patchcmd{\l@section}{#1}{\MakeUppercase{#1}}{}{}
\makeatother
\renewcommand{\cftsecleader}{\cftdotfill{\cftdotsep}}
\renewcommand{\cftsubsecleader}{\cftdotfill{\cftdotsep}}
\renewcommand{\cftsubsubsecleader}{\cftdotfill{\cftdotsep}}
\renewcommand{\cftsecfont}{\fontsize{12}{14}\bfseries}
\renewcommand{\cftsecpagefont}{\fontsize{12}{14}\bfseries}
\renewcommand{\cftsecaftersnum}{.\enspace}
\settowidth{\cftsecnumwidth}{VIII.\enspace}
\renewcommand{\cftsubsecaftersnum}{.}
\renewcommand{\cftsubsubsecaftersnum}{.}

% --- Надписи: таблици ОТГОРЕ, фигури ОТДОЛУ (А.8, А.9) ------------------
\usepackage{caption}
\usepackage{subcaption}
\DeclareCaptionLabelFormat{tablelabel}{Таблица #2}
\captionsetup[table]{labelformat=tablelabel, labelsep=period, position=above,
                     font={bf,small}, justification=centering, singlelinecheck=true}
\DeclareCaptionLabelFormat{figurelabel}{Фигура #2}
\captionsetup[figure]{labelformat=figurelabel, labelsep=period, position=below,
                      font={bf,small}, justification=centering, singlelinecheck=true}

% Препратки в текста: „(Табл. 3)“ / „(Фиг. 7)“, задължителни по А.8/А.9.
\newcommand{\tabref}[1]{(Табл.~\ref{#1})}
\newcommand{\figref}[1]{(Фиг.~\ref{#1})}

% --- Математика ---------------------------------------------------------
\usepackage{amsmath,amssymb,amsfonts}
\usepackage{amsthm}

% --- Графики ------------------------------------------------------------
\usepackage{graphicx}
\usepackage{float}
\usepackage{tikz}
\usetikzlibrary{decorations.pathreplacing,arrows.meta,positioning,shapes.geometric,fit,calc}
\usepackage{adjustbox}
\usepackage{pdfpages}     % за подвързаните задание и декларация, вж. Б.5

% --- Таблици ------------------------------------------------------------
\usepackage{array}
% Ляво подравнена колона с фиксирана ширина. Тесните `p{}` колони с ДВУСТРАННО
% подравняване (наследено от основния текст) разкъсват редовете с огромни
% междудумни интервали. Правилото за justified от А.4 се отнася за ОСНОВНИЯ текст,
% не за клетки от 3 cm. Всички таблици с текстови клетки ползват L.
\newcolumntype{L}[1]{>{\RaggedRight\arraybackslash}p{#1}}
\usepackage{booktabs}
\usepackage[table]{xcolor}% \cellcolor за топлинната карта
\usepackage{longtable}    % „продължение“ при пренасяне на таблица (А.8)
\usepackage{multirow}
\usepackage{tabularx}

% --- Помощни ------------------------------------------------------------
\usepackage{hyperref}
\usepackage{xurl}
\usepackage{enumitem}
\usepackage{textcomp}
\usepackage{listings}
\urlstyle{same}
% КРИТИЧНО: `breaklines` на listings вмъква ДИСКРЕЦИОННИ прекъсвания, а глобалните
% \hyphenpenalty/\exhyphenpenalty=10000 (правилото „без пренасяния“, А.4) ги правят
% безкрайно скъпи, тогава дългите редове в листингите НЕ се пренасят и излизат извън
% полето. Проверено изолирано: с наказанията 1 препълване, без тях 0. Наказанията се
% нулират САМО вътре в листинг, така че текстът си остава без пренасяния.
\makeatletter
\lst@AddToHook{PreInit}{\hyphenpenalty=50\exhyphenpenalty=50}
\makeatother
\renewcommand{\lstlistingname}{Листинг}
\renewcommand{\lstlistlistingname}{Списък на листингите}

% Italic САМО за ключови понятия (А.4), семантичен макрос, за да е ревизуемо.
\newcommand{\term}[1]{\emph{#1}}
% Инлайн код/идентификатори.
\DeclareRobustCommand{\code}[1]{\texttt{#1}}
% Глифи за таблиците с присъди (генерирани от tools/gen_tex.py).
\newcommand{\cmark}{\ensuremath{\checkmark}}
\newcommand{\xmark}{\ensuremath{\times}}

% Слот за генерирани макроси с числа (за да не може изречение да противоречи на
% таблица). Празен е, докато измервателният скрипт не започне да пише в generated/.
% \input{generated/metrics_macros}

% Заснемане с тънка черна рамка. Рамката НЕ е украса: интерфейсите със светла тема са
% почти бели и се сливат с бялата страница, така че без нея не се вижда къде свършва
% изображението и започва листът, а размерът на изображението е част от твърдението.
% \fboxsep=0pt прилепва рамката към пиксела; 0.4 pt е най-тънката, която се възпроизвежда
% надеждно при печат.
\newcommand{\capture}[2][0.44\textwidth]{%
  {\setlength{\fboxsep}{0pt}\setlength{\fboxrule}{0.4pt}%
   \fbox{\includegraphics[width=#1]{#2}}}}

% Съпоставка „изходен код | диаграма“ за един пример.
%   #1 папка на примера  #2 ключ на примера  #3 описание
\newcommand{\roundtrippair}[3]{%
  \begin{figure}[H]\centering
    \begin{tabular}{@{}cc@{}}
      \small\textbf{изходен код} & \small\textbf{диаграма} \\
      \capture{figures/#1/#2_code.png} &
      \capture{figures/#1/#2_uml.png} \\
    \end{tabular}
    \caption{#3}\label{fig:rt-#1-#2}
  \end{figure}}

% Видим маркер за незапълнено място. Търси се с: grep -rn 'TODO' chapters/
\newcommand{\TODO}[1]{\textcolor{red}{\textbf{[TODO: #1]}}}

\tikzset{
    block/.style={draw, rounded corners, fill=blue!5, align=center, font=\small,
        minimum width=3.2cm, minimum height=1.05cm, inner sep=6pt},
    accent/.style={draw, rounded corners, fill=green!8, align=center, font=\small,
        minimum width=3.0cm, minimum height=0.95cm, inner sep=6pt},
    danger/.style={draw, rounded corners, fill=red!6, align=center, font=\small,
        minimum width=3.0cm, minimum height=0.95cm, inner sep=6pt},
    note/.style={draw, dashed, rounded corners, fill=gray!8, align=left, font=\small,
        inner sep=6pt},
    line/.style={-Latex, thick}
}

% --- Единично разстояние в нетекстови среди -----------------------------
\AtBeginEnvironment{tikzpicture}{\setstretch{1.1}}
\AtBeginEnvironment{lstlisting}{\setstretch{1.1}}
\AtBeginEnvironment{verbatim}{\setstretch{1.1}}
\AtBeginEnvironment{tabular}{\setstretch{1.1}}
\AtBeginEnvironment{longtable}{\setstretch{1.1}}

\BeforeBeginEnvironment{tikzpicture}{\begin{adjustbox}{max width=\textwidth}}
\AfterEndEnvironment{tikzpicture}{\end{adjustbox}}

% --- Оцветяване на C++ код ----------------------------------------------
\definecolor{cppblue}{RGB}{0,0,180}
\definecolor{cppgreen}{RGB}{0,128,0}
\definecolor{cppred}{RGB}{180,0,0}
\definecolor{cppgray}{RGB}{128,128,128}

\lstset{
    language=C++,
    basicstyle=\footnotesize\ttfamily,
    % Пренасяне на дълги редове. `breakatwhitespace=false` е СЪЩЕСТВЕНО: коментарите в
    % изходния текст стигат 110 знака без удобно място за прекъсване, и при `true` редът
    % изобщо не се пренася, а излиза извън полето. `linewidth` се задава изрично, защото
    % рамката и левият отстъп иначе не се приспадат от наличната ширина.
    breaklines=true, breakatwhitespace=false, breakautoindent=false, breakindent=0pt,
    linewidth=\dimexpr\linewidth-2.5em\relax,
    columns=flexible, keepspaces=true,
    postbreak=\mbox{\textcolor{cppgray}{\ensuremath{\hookrightarrow}}\space},
    xleftmargin=1.5em, xrightmargin=0.5em,
    keywordstyle=\color{cppblue}\bfseries,
    stringstyle=\color{cppred},
    commentstyle=\color{cppgreen},
    identifierstyle=\color{black},
    showstringspaces=false,
    numbers=left, numberstyle=\tiny\color{cppgray}, numbersep=8pt,
    aboveskip=1.5em, belowskip=1.5em,
    frame=single, framerule=0.5pt, rulecolor=\color{cppgray},
    tabsize=4, captionpos=b,
    morekeywords={concept, requires, constexpr, consteval, constinit, noexcept,
                  nullptr, override, final, co_await, co_return, co_yield,
                  expected, unexpected, span, string_view, unique_ptr, shared_ptr},
    emph={blueprint, model, class_node, attribute, operation, visibility, relation,
          association, aggregation, composition, generalization, stereotype,
          emitter, plantuml_emitter, mermaid_emitter, schema, table, column,
          foreign_key, consistency_checker, cursor, translation_unit},
    emphstyle=\color{cppblue},
    inputencoding=utf8,
    % Не-ASCII знаци, срещащи се в коментарите на изходния текст. Без тези редове
    % \lstinputlisting отказва с „Invalid UTF-8 byte sequence“, listings не разбира
    % многобайтов вход, който не е в literate таблицата.
    literate={→}{{$\rightarrow$}}1 {×}{{$\times$}}1 {…}{{...}}1 {§}{{\S}}1
{а}{{\cyra}}1 {б}{{\cyrb}}1 {в}{{\cyrv}}1 {г}{{\cyrg}}1 {д}{{\cyrd}}1
             {е}{{\cyre}}1 {ж}{{\cyrzh}}1 {з}{{\cyrz}}1 {и}{{\cyri}}1 {й}{{\cyrishrt}}1
             {к}{{\cyrk}}1 {л}{{\cyrl}}1 {м}{{\cyrm}}1 {н}{{\cyrn}}1 {о}{{\cyro}}1
             {п}{{\cyrp}}1 {р}{{\cyrr}}1 {с}{{\cyrs}}1 {т}{{\cyrt}}1 {у}{{\cyru}}1
             {ф}{{\cyrf}}1 {х}{{\cyrh}}1 {ц}{{\cyrc}}1 {ч}{{\cyrch}}1 {ш}{{\cyrsh}}1
             {щ}{{\cyrshch}}1 {ъ}{{\cyrhrdsn}}1 {ь}{{\cyrsftsn}}1 {ю}{{\cyryu}}1 {я}{{\cyrya}}1
             {А}{{\CYRA}}1 {Б}{{\CYRB}}1 {В}{{\CYRV}}1 {Г}{{\CYRG}}1 {Д}{{\CYRD}}1
             {Е}{{\CYRE}}1 {Ж}{{\CYRZH}}1 {З}{{\CYRZ}}1 {И}{{\CYRI}}1 {Й}{{\CYRISHRT}}1
             {К}{{\CYRK}}1 {Л}{{\CYRL}}1 {М}{{\CYRM}}1 {Н}{{\CYRN}}1 {О}{{\CYRO}}1
             {П}{{\CYRP}}1 {Р}{{\CYRR}}1 {С}{{\CYRS}}1 {Т}{{\CYRT}}1 {У}{{\CYRU}}1
             {Ф}{{\CYRF}}1 {Х}{{\CYRH}}1 {Ц}{{\CYRC}}1 {Ч}{{\CYRCH}}1 {Ш}{{\CYRSH}}1
             {Щ}{{\CYRSHCH}}1 {Ъ}{{\CYRHRDSN}}1 {Ь}{{\CYRSFTSN}}1 {Ю}{{\CYRYU}}1 {Я}{{\CYRYA}}1,
}


\begin{document}

% ============================================================================
% Заглавна страница (А.5)
% ============================================================================
\begin{titlepage}
\begin{center}
    {\fontsize{18}{21.6}\selectfont\bfseries ТЕХНИЧЕСКИ УНИВЕРСИТЕТ \textendash{} СОФИЯ\par}
    \vspace{0.5cm}
    {\fontsize{16}{19.2}\selectfont\bfseries ФАКУЛТЕТ „КОМПЮТЪРНИ СИСТЕМИ И ТЕХНОЛОГИИ“\par}

    \vspace{0.3cm}
    \begin{figure}[H]
        \centering
        \includegraphics[width=0.30\linewidth]{images/tu-logo.png}
    \end{figure}
    \vspace{0.2cm}

    {\fontsize{20}{23}\selectfont\bfseries ДВУПОСОЧНО ИНЖЕНЕРСТВО МЕЖДУ ИЗХОДЕН КОД НА C++, UML ДИАГРАМИ НА КЛАСОВЕ И РЕЛАЦИОННИ СХЕМИ НА ДАННИ\par}

    \vspace{0.6cm}

    % NB: формулировката може да се наложи да бъде „КУРСОВА РАБОТА“ вместо
    % „КУРСОВ ПРОЕКТ“. Учебният план изброява и двете форми и решаващото е
    % какво възлага самата дисциплина. Да се сверява преди печат.
    {\fontsize{14}{16.8}\selectfont\bfseries КУРСОВ ПРОЕКТ\par}
    \vspace{0.3cm}
    {\fontsize{14}{16.8}\selectfont\bfseries по дисциплина „UML обектно-ориентирано проектиране“\par}

    \vfill

    \begin{flushleft}
    {\fontsize{14}{16.8}\selectfont\bfseries Разработил:\par}
    Алекс Цветанов, фак. № \TODO{фак. номер}\par
    Специалност: Компютърно и софтуерно инженерство, ОКС „Магистър“\par

    \vspace{0.7cm}

    {\fontsize{14}{16.8}\selectfont\bfseries Преподавател:\par}
    Проф. д-р инж. Даниела Гоцева\par
    \end{flushleft}

    \vfill

    % Годината на предаване. Да се приведе към действителната дата на защита
    % на курсовия проект.
    {\fontsize{12}{14.4}\selectfont София, \TODO{година на предаване}\par}
\end{center}
\end{titlepage}

% ============================================================================
% Съдържание (А.3, А.6)
% ============================================================================
\newpage
\pagestyle{empty}
\setcounter{tocdepth}{2}
\tableofcontents
\newpage

\pagestyle{fancy}
\pagenumbering{arabic}

% ============================================================================
% УВОД, без номер
% ============================================================================
% УВОД, без номер. Ориентировъчен обем: 1 до 2 стр.
% Резюме на дейностите по курсовия проект: задача, цел, задачи, обхват, структура.
\section*{УВОД}
\label{sec:introduction}
\addcontentsline{toc}{section}{УВОД}

Обектно-ориентираният модел и програмният код, който го реализира, живеят в два
различни носителя. Моделът се описва на UML и служи за разговор между хора, а кодът се
описва на език за програмиране и служи за разговор с машина. Двата носителя се
разминават почти веднага след като работата по проекта започне, защото кодът се променя
всеки ден, а диаграмата се преначертава, когато някой се сети. Резултатът е познат:
диаграма, на която вече никой не вярва, и код, чиято структура се възстановява само с
четене.

Настоящият курсов проект разглежда това разминаване като инженерна задача с решение, а
не като организационен навик. Ако структурата на кода може да бъде извлечена машинно, а
UML моделът може да бъде записан в текстов вид, то преходът между двете посоки е
преобразувание между два формални езика и подлежи на автоматизация. Терминологичната
рамка за двете посоки е утвърдена от Chikofsky и Cross~\cite{chikofsky1990reverse}:
\term{обратно инженерство} е възстановяване на абстракция от реализация, а
\term{право инженерство} е обратният преход от абстракция към реализация. Двете заедно
образуват \term{двупосочно инженерство}, което е предметът на този проект.

\textbf{Целта} на проекта е да се проектира и реализира инструмент, наречен Blueprint,
който извлича UML модел на класове от изходен код на C++, изобразява го в текстови
нотации за диаграми, извежда от него релационна схема на данни и в обратната посока
поражда заготовки на класове и схема по подаден модел, като допълнително отчита
местата, на които моделът и кодът са се разминали.

За постигането на целта са формулирани следните \textbf{задачи}:

\begin{enumerate}[topsep=0pt,itemsep=0pt,leftmargin=*]
\item преглед на понятийния модел на UML, на градивните му елементи, на видовете
      връзки и на правилата и общите механизми, доколкото те засягат диаграмата на
      класове и нейното съответствие с изходен код;
\item анализ на съществуващите инструменти за извличане на структура от код и за
      двупосочно инженерство, и обосновка на избора на подход и на технологичен набор;
\item проектиране на вътрешния модел на системата, на границите между модулите и на
      правилата за преобразуване между клас, диаграма и релационна схема;
\item програмна реализация на извличането, на извеждането в двете текстови нотации, на
      пораждането на схема и заготовки и на проверката за разминаване;
\item изграждане на методика за експериментална проверка на верността на
      преобразуванията и на устойчивостта им при преминаване през пълен цикъл;
\item провеждане на експериментите и анализ на получените резултати.
\end{enumerate}

\textbf{Обхватът} на работата е структурната част на обектно-ориентирания модел:
класове, атрибути, операции, видимост, обхват, шаблонни класове и връзките между
класовете, заедно със съответствието им на релационна схема. Диаграмите на случаите на
употреба и на дейностите се разглеждат като контекст на модела и се използват за
описание на самата система, но не се извличат от изходен код. Извън обхвата остават
поведенческите диаграми, които изискват анализ на потока на управление, както и пълното
семантично покритие на езика C++.

\textbf{Структура на изложението.} Раздел~\ref{sec:literature} излага предметната област
и понятийния апарат на UML. Раздел~\ref{sec:alternatives} съпоставя възможните решения и
обосновава избора. Раздел~\ref{sec:design} представя проектирането на системата, а
Раздел~\ref{sec:implementation} нейната програмна реализация.
Раздел~\ref{sec:method} описва методиката за проверка, а Раздел~\ref{sec:results}
получените резултати. Заключението обобщава постигнатото и очертава по-нататъшната
работа.

\newpage

% ============================================================================
% I. Предметна област и обосновка на избора (слят: преглед и алтернативи)
% ============================================================================
% I. ПРЕДМЕТНА ОБЛАСТ И ЛИТЕРАТУРЕН ПРЕГЛЕД
% Покрива от учебната програма: характеристики на UML, понятиен модел, градивни
% елементи, връзки, диаграми, правила и общи механизми, моделиране на системи,
% жизнен цикъл на програмния продукт и UML.
\section{Предметна област и литературен преглед}
\label{sec:literature}

UML е графичен език за визуализиране, специфициране, конструиране и документиране на
артефактите на една програмна система. Той не е метод и не предписва процес на
разработка, а речник и правила за съчетаване на думите от този речник. Тази разлика е
съществена за настоящия проект: щом езикът е формален, а не илюстративен, той има
недвусмислена абстрактна структура и следователно подлежи на машинна обработка.
Официалната спецификация на езика~\cite{omg_uml251} задава именно тази структура под
формата на метамодел, а ръководствата на авторите на
езика~\cite{booch2005userguide,rumbaugh2004reference} го излагат в приложим за
проектант вид.

Понятийният модел на UML стъпва на три групи съставки. Първата са \term{градивните
елементи}: структурни (клас, интерфейс, компонент, възел), поведенчески
(взаимодействие, състояние) и групиращи (пакет), заедно с пояснителния елемент
(бележка). Втората са \term{връзките}: зависимост, асоциация, обобщение и реализация,
като асоциацията има отделни разновидности за агрегация и композиция, различаващи се по
това дали частта може да съществува без цялото. Третата са \term{правилата и общите
механизми}: спецификации, украшения, общи разделения и \term{механизмите за
разширяване}, тоест ограничения, етикетирани стойности и стереотипи. Настоящият проект
се интересува непосредствено от последните, защото именно чрез стереотип и етикетирана
стойност се обозначава кой клас участва в релационната схема и как.

Диаграмата на класове е основната структурна диаграма и е носител на речника на
системата. Класификаторът в нея се описва с име, атрибути, операции и отговорности,
като видимостта и обхватът се отбелязват върху всеки член. Практиката за откриване на
този речник е по-стара от самия UML: картите клас, отговорност и сътрудничество, познати
като CRC карти, са предложени от Beck и Cunningham~\cite{beck1989crc} точно като начин
класовете да бъдат намерени преди да бъдат нарисувани. За компактно въведение в
приложната част на езика е използван и~\cite{fowler2003distilled}.

Отношението между обектния модел и релационния модел на данните е самостоятелна тема със
собствена литература. Релационната схема се описва чрез таблици, ключове и връзки между
тях и има свой понятиен предшественик в модела същност и връзка на
Chen~\cite{chen1976er}. Съответствието между клас и таблица не е взаимно еднозначно:
наследяването няма пряк аналог в релационния модел, а асоциацията се изразява чрез външен
ключ или чрез свързваща таблица според кратността. Систематичен преглед на утвърдените
стратегии за това съответствие дава~\cite{ambler2003agiledb}. Тези стратегии са
непосредствената основа на правото инженерство в настоящия проект.

Мястото на двупосочното инженерство в жизнения цикъл на програмния продукт следва
таксономията на Chikofsky и Cross~\cite{chikofsky1990reverse}. Обратното инженерство
поддържа документацията вярна, без да изисква дисциплина от разработчика, и е приложимо
и към чужд код. Правото инженерство пренася решение, взето на нивото на модела, надолу
до реализацията, без ръчно преписване. Стойността на двете заедно е в третото им
следствие: щом и двете посоки са автоматизирани, разминаването между модел и код става
измеримо, а не просто подозирано. \TODO{при разширяване на раздела да се добави
преглед на инструментите за проверка на съответствие между модел и код}

\newpage

% ============================================================================
% II. Проектиране и програмна реализация.
% Двата файла са ЕДИН раздел: 03 дава решенията, 04 дава реализацията им.
% Само 03 отваря \section, 04 продължава без нов заглавен ред и без \newpage.
% ============================================================================
% III. ПРОЕКТИРАНЕ НА СИСТЕМАТА
% Тук се събира по-голямата част от покритието на учебната програма: класове,
% операции, атрибути, отговорности, CRC карти, речник на системата, връзки,
% зависимост и наследяване, ограничения, етикетирани стойности и стереотипи,
% класификатори, видимост и обхват, шаблонни класове, диаграми на случаите на
% употреба и на дейностите.
\section{Проектиране и програмна реализация}
\label{sec:design}

\subsection{Модули и посока на зависимостите}

Системата е пет модула с еднопосочни зависимости. \term{Четецът} обхожда единицата за превод и
произвежда вътрешен модел. \term{Моделът} е независим от езика на входа и от формата на изхода и е
единственото място със структура. \term{Излъчвателите} го превръщат в текст на избраната нотация,
\term{пораждащият модул} в релационна схема и в заготовки, а \term{проверяващият модул} съпоставя два
модела. Зависимостите сочат само към модела, така че нито един излъчвател не знае за съществуването
на четеца \figref{fig:packages}.

\begin{figure}[H]
\centering
\begin{tikzpicture}[node distance=1.0cm and 1.2cm]
  \node[block] (model) {\textbf{модел}\\ класификатори, членове,\\ връзки};
  \node[accent, left=2.4cm of model] (reader) {\textbf{четец}\\ AST в JSON\\ към модел};
  \node[accent, above=1.1cm of model] (diagramreader) {\textbf{четец на диаграми}\\ текст към модел};
  \node[accent, right=2.4cm of model] (emit) {\textbf{излъчватели}\\ PlantUML, Mermaid};
  \node[accent, below left=1.2cm and 0.2cm of model] (gen) {\textbf{пораждащ модул}\\ схема и заготовки};
  \node[accent, below right=1.2cm and 0.2cm of model] (check) {\textbf{проверяващ модул}\\ разлики};
  \node[note, left=1.0cm of reader, text width=2.1cm] (src) {единица\\ за превод};
  \node[note, right=1.0cm of emit, text width=2.1cm] (out) {текст на\\ диаграма};

  \draw[line] (src) -- (reader);
  \draw[line] (reader) -- (model);
  \draw[line] (diagramreader) -- (model);
  \draw[line] (model) -- (emit);
  \draw[line] (emit) -- (out);
  \draw[line] (model) -- (gen);
  \draw[line] (model) -- (check);
\end{tikzpicture}
\caption{Модули и посока на зависимостите. Всяка стрелка сочи натам, накъдето тече информация, и нито един модул не зависи от друг освен от модела}
\label{fig:packages}
\end{figure}

Правилото се проверява механично, а не се пази по споразумение: заглавните файлове на излъчвателите,
на пораждащия и на проверяващия модул включват \code{model.hpp} и не включват \code{reader.hpp}.

\subsection{Речник, класификатори и връзки}

Речникът е получен по метода на CRC картите~\cite{beck1989crc}: за всяка предполагаема отговорност е
записан носителят ѝ и сътрудниците му, след което картите без собствено поведение са отпаднали или са
се слели със съседите си. Класификаторът носи квалифицирано име, вид, шаблонни параметри, атрибути,
операции, стереотипи и етикетирани стойности; атрибутът име, тип, видимост, обхват и кратност;
операцията и признаци за виртуалност, чистота и константност. Статичните членове се маркират като
членове с обхват на класификатора. Видовете класификатори са пет: клас, структура, интерфейс, изброен
тип и шаблон. Интерфейсът не се чете от изходния текст, защото C++ няма такава ключова дума, а се
\term{извежда}: класификатор без атрибути, чиито операции са само чисти виртуални, е интерфейс по
определението на UML, и изводът се прави върху модела, а не върху синтаксиса.

Връзките се извеждат от типовете на членовете, а не се обявяват отделно. Обобщението следва пряко
от списъка на базовите класове, заедно с достъпа на всяка база и с признака за виртуално
наследяване. Зависимостта идва от употребата на тип в подпис на операция, когато няма по-силна
връзка, а реализацията свързва специализация с общия шаблон. Разграничението между асоциация,
агрегация и композиция се прави по притежанието и е изнесено в таблица, за да бъде ревизуемо
поотделно \tabref{tab:ownership}.

\begin{table}[H]
\caption{Съответствие между изписването на члена и вида на връзката}
\label{tab:ownership}
\centering
\fontsize{10}{11}\selectfont
\setstretch{1.0}
\begin{tabular}{@{}L{4.6cm}L{2.6cm}L{3.0cm}L{1.8cm}@{}}
\toprule
\textbf{Изписване на члена} & \textbf{Притежание} & \textbf{Връзка} & \textbf{Кратност} \\
\midrule
\code{T value}                    & изключително & композиция  & 1 \\
\code{std::unique\_ptr<T>}        & изключително & композиция  & 0..1 \\
\code{std::optional<T>}           & изключително & композиция  & 0..1 \\
\code{std::shared\_ptr<T>}        & споделено    & агрегация   & 0..1 \\
\code{std::weak\_ptr<T>}          & споделено    & агрегация   & 0..1 \\
\code{T*}                         & няма         & асоциация   & 0..1 \\
\code{T\&}                        & няма         & асоциация   & 1 \\
контейнер от някое от горните     & като елемента & като елемента & * \\
тип само в подпис на операция     & няма         & зависимост  & \\
\bottomrule
\end{tabular}
\end{table}

Прочитът следва UML буквално. Композиция означава, че частта загива заедно с цялото, затова член по
стойност или в указател с изключително притежание е композиция; \code{shared\_ptr} задържа частта
жива и след цялото, което е агрегация; обикновен указател или отправка не твърдят нищо за времето на
живот, което оставя проста асоциация. Член, чийто тип не е класификатор от модела, остава обикновен
атрибут и не поражда връзка. Шаблонните класове се записват като параметризирани класификатори, а
специализациите като отделни класификатори със връзка на реализация към общия шаблон. Разгръщането
на всяка инстанциация в отделен клас би затрупало диаграмата с еднотипни възли и би заличило самата
параметризация, която е носителят на смисъла; по същата причина \term{инстанциациите, породени от
употребата}, не влизат в модела, тъй като са свойство на конкретния превод, а не на проектното
решение.

\subsection{Съответствие с релационна схема и поведенчески модели}

Клас със стереотипа за постоянно съхранение поражда таблица, всеки негов атрибут със скаларен тип
поражда колона, а етикетираните стойности задават име на таблицата, име и тип на колоната и участие в
първичния ключ. Клас без обозначен първичен ключ получава сурогатен ключ \code{id}, и това се съобщава
изрично. Асоциацията с кратност едно поражда външен ключ в притежаващата страна, а с кратност много
свързваща таблица; изборът е съзнателен, защото диаграмата показва само единия край, тоест от нея не
се вижда дали връзката е едно към много или много към много, а свързващата таблица е вярна и в двата
случая.

За наследяването са реализирани и трите утвърдени стратегии~\cite{ambler2003agiledb}: една таблица за
цялата йерархия с колона разграничител, таблица на всеки клас с ключ, който е и външен ключ към базата,
и таблица на всеки конкретен клас с копирани надолу колони. По подразбиране се прилага \term{една
таблица за йерархията}, защото само при нея извличането на един обект не изисква съединяване; изборът
не е зазидан, а се прави с етикетирана стойност върху базовия клас.

Нито диаграмата на случаите на употреба \figref{fig:usecase}, нито тази на дейностите
\figref{fig:activity} може да бъде извлечена от изходен текст: случаят на употреба е твърдение за
намерение, а дейността за поток на управление, който структурата на кода не носи. Затова двете се
описват на ръка и се превеждат към същите две нотации като диаграмата на класове.

\begin{figure}[htbp]
\centering
\begin{adjustbox}{max width=\linewidth, max totalheight=0.36\textheight}
\begin{tikzpicture}[node distance=0.7cm and 2.0cm]
  \tikzset{bpuc/.style={draw, ellipse, align=center, font=\normalsize, minimum width=5.4cm,
                      minimum height=1.05cm, fill=blue!5, inner sep=2pt},
           bpact/.style={draw, align=center, font=\normalsize, minimum width=2.1cm,
                       minimum height=0.8cm, fill=green!8, rounded corners}}
  \node[bpuc] (uc1) {UC1: извличане на модел\\ от изходен текст};
  \node[bpuc, below=of uc1] (uc2) {UC2: излъчване на\\ диаграма на класове};
  \node[bpuc, below=of uc2] (uc3) {UC3: пораждане на\\ релационна схема};
  \node[bpuc, below=of uc3] (uc4) {UC4: пораждане на заготовки\\ по схема};
  \node[bpuc, below=of uc4] (uc5) {UC5: прочитане на диаграма\\ обратно в модел};
  \node[bpuc, below=of uc5] (uc6) {UC6: отчитане на разминаване\\ между диаграма и код};
  \node[bpuc, below=of uc6] (uc7) {UC7: избор на стратегия\\ за наследяването};

  \node[bpact, left=of uc2] (dev) {Разработчик};
  \node[bpact, left=of uc5] (arch) {Проектант};
  \node[bpact, below=1.4cm of arch] (ci) {Конвейер за\\ интеграция};

  \node[draw, dashed, rounded corners, fit=(uc1)(uc7), inner sep=9pt,
        label={[font=\normalsize\bfseries]above:Blueprint}] (sys) {};

  \foreach \a/\b in {dev/uc1, dev/uc2, dev/uc3, dev/uc4, arch/uc2, arch/uc5, arch/uc6, ci/uc6}
      \draw (\a) -- (\b);

  % Отношенията между случаите се извеждат вдясно от рамката на системата,
  % за да не минават през самите елипси.
  \draw[line, dashed] (uc2.east) -- ++(0.7,0)
        node[right, font=\small, xshift=1pt] {$\ll$include$\gg$} |- (uc1.east);
  \draw[line, dashed] (uc3.east) -- ++(2.6,0)
        node[above, font=\small, yshift=2pt] {$\ll$include$\gg$} |- (uc1.east);
  \draw[line, dashed] (uc6.east) -- ++(0.7,0)
        node[right, font=\small, xshift=1pt] {$\ll$include$\gg$} |- (uc5.east);
  \draw[line, dashed] (uc7.east) -- ++(4.5,0)
        node[above, font=\small, yshift=2pt] {$\ll$extend$\gg$} |- (uc3.east);
\end{tikzpicture}
\end{adjustbox}
\caption{Случаи на употреба. Изходният текст на модела е в \code{examples/blueprint.usecase} и се превежда към двете нотации с \code{blueprint usecase}}
\label{fig:usecase}
\end{figure}

\begin{figure}[htbp]
\centering
\begin{adjustbox}{max width=\linewidth, max totalheight=0.36\textheight}
\begin{tikzpicture}[node distance=0.75cm and 1.1cm]
  \tikzset{bpstep/.style={draw, rounded corners, align=center, font=\normalsize,
                        minimum width=5.2cm, minimum height=0.85cm, fill=blue!5},
           bpdec/.style={draw, diamond, aspect=2.6, align=center, font=\small,
                       fill=yellow!12, inner sep=1pt},
           bpdot/.style={draw, circle, fill=black, minimum size=5pt, inner sep=0pt}}
  \node[bpdot] (s) {};
  \node[bpstep, below=of s] (a1) {Изпълнение на фронтенда върху\\ единицата за превод};
  \node[bpstep, below=of a1] (a2) {Прочитане на синтактичното дърво\\ в модела};
  \node[bpdec, below=of a2] (d1) {има ли класификатор\\ със стереотип за съхранение};
  \node[bpstep, below left=1.0cm and 0.4cm of d1] (a3) {Излъчване на диаграмата\\ в двете нотации};
  \node[bpstep, below right=1.0cm and 0.4cm of d1] (a4) {Извеждане на\\ релационната схема};
  \node[bpstep, below=of a3] (a7) {Прочитане на диаграмата\\ обратно в модел};
  \node[bpstep, below=of a4] (a5) {Пораждане на заготовки\\ по схемата};
  \node[bpstep, below=of a5] (a6) {Извличане на модел от\\ породените заготовки};
  \node[bpstep, below=1.6cm of a7] (a8) {Съпоставяне на двата модела};
  \node[bpdec, below=of a8] (d2) {съвпадат ли};
  \node[bpstep, below left=0.9cm and 0.1cm of d2] (a10) {Отчитане на затворен цикъл};
  \node[bpstep, below right=0.9cm and 0.1cm of d2] (a9) {Отчитане на разминаване\\ и прекъсване на построяването};
  \node[bpdot, below=1.1cm of d2, yshift=-1.0cm, minimum size=6pt, label={[bpdot, minimum size=11pt, fill=none]center:}] (e) {};

  \draw[line] (s) -- (a1);
  \draw[line] (a1) -- (a2);
  \draw[line] (a2) -- (d1);
  \draw[line] (d1) -| node[above left, font=\small] {не} (a3);
  \draw[line] (d1) -| node[above right, font=\small] {да} (a4);
  \draw[line] (a3) -- (a7);
  \draw[line] (a4) -- (a5);
  \draw[line] (a5) -- (a6);
  \draw[line] (a7) -- (a8);
  \draw[line] (a6) |- (a8);
  \draw[line] (a8) -- (d2);
  \draw[line] (d2) -| node[above left, font=\small] {да} (a10);
  \draw[line] (d2) -| node[above right, font=\small] {не} (a9);
  \draw[line] (a10) |- (e);
  \draw[line] (a9) |- (e);
\end{tikzpicture}
\end{adjustbox}
\caption{Дейности при преминаване през пълен цикъл. Изходният текст на модела е в \code{examples/roundtrip.activity} и се превежда с \code{blueprint activity}}
\label{fig:activity}
\end{figure}

% IV. ПРОГРАМНА РЕАЛИЗАЦИЯ
% Обща структура, йерархия на класовете, коментар на изходния текст (по А.3).
% Листингите се вмъкват с \lstinputlisting от src/, за да не се преписват.
\section{Програмна реализация}
\label{sec:implementation}

\subsection{Обща структура}

Реализацията е на C++20~\cite{iso_cpp20} и се построява с CMake. Проектът е разделен на
библиотека, която съдържа цялата логика, и изпълним файл, който е тънък слой над нея за
работа от команден ред. Разделението не е украса: модулните тестове се свързват срещу
библиотеката и не минават през командния ред, така че проверката на поведението е
независима от начина на извикване. Публичните заглавни файлове живеят в \code{include/%
blueprint/}, реализацията в \code{src/}, а тестовете в \code{tests/}. Всеки модул от
Раздел~\ref{sec:design} има точно един заглавен файл и точно един файл с реализация
\tabref{tab:files}.

\begin{table}[H]
\caption{Файлове на реализацията и отговорността на всеки}
\label{tab:files}
\centering
\fontsize{11}{13}\selectfont
\begin{tabular}{@{}L{3.9cm}L{8.1cm}@{}}
\toprule
\textbf{Файл} & \textbf{Отговорност} \\
\midrule
\code{json.\{hpp,cpp\}} & Документен модел, четец и записвач на JSON. Съществува, защото изводът на фронтенда е JSON, а проектът не приема външни зависимости. \\
\code{model.\{hpp,cpp\}} & Вътрешният модел и записът му във външен вид. Единственото място, където се пази структура. \\
\code{reader.\{hpp,cpp\}} & Обхождане на синтактичното дърво, извеждане на връзките, изпълнение на компилатора. \\
\code{emit.\{hpp,cpp\}} & Общият обход и двата излъчвателя. \\
\code{diagram.\{hpp,cpp\}} & Прочитане на диаграма обратно в модел, за двете нотации. \\
\code{schema.\{hpp,cpp\}} & Модел към релационна схема, схема към DDL и обратно, схема към заготовки на класове. \\
\code{diff.\{hpp,cpp\}} & Съпоставяне на два модела и отчитане на разликите. \\
\code{behavior.\{hpp,cpp\}} & Модели и излъчватели за случаите на употреба и за дейностите. \\
\code{annotations.hpp} & Макросите за анотиране, само дефиниции на препроцесора. \\
\code{io.\{hpp,cpp\}} & Четене и запис на файл, за да се отказват еднакво навсякъде. \\
\code{main.cpp} & Команден ред. Не взема решения, само подрежда извиквания. \\
\bottomrule
\end{tabular}
\end{table}

\subsection{Четене на изходния текст}

Четенето става чрез изпълнение на компилатора като външен процес и разбор на изведеното
от него синтактично дърво в JSON~\cite{clang_ast_json}. Съставянето на командата е
отделено в \code{build\_clang\_command}, което я прави проверима без да се изпълнява
процес, а разборът на резултата е отделен в \code{read\_ast\_json}, което е чиста функция
без достъп до файловата система. Разделянето не е формално: то позволява целият четец да
се тества върху записан образец на дърво, без инсталиран компилатор.

Три подробности от обхождането заслужават коментар, защото и трите са места, на които
наивната реализация мълчаливо произвежда грешен модел.

\textbf{Проследяване на текущия файл.} Фронтендът извежда полето с името на файла
\term{само когато то се променя}, а следващите възли го наследяват. Затова текущият файл
се пренася през цялото обхождане в реда на документа. Пропускането на тази подробност е
разликата между модел на проекта и модел на всеки системен заглавен файл, който проектът
включва.

\textbf{Отсяване на подразбиращите се членове.} Фронтендът извежда и членовете, които
самият той е породил: подразбиращите се конструктори, операторите за присвояване,
деструкторът, както и вграденото име на класа в самия него. Всички те носят признака
\code{isImplicit} и се отхвърлят. Без това всеки клас в диаграмата получава по пет
операции, които никой не е писал.

\textbf{Видимост по секции.} Видимостта на един член не се извежда като негово поле, а
следва от предхождащия го спецификатор за достъп. Затова тя се води при обхождането,
започвайки от \code{private} за \code{class} и от \code{public} за \code{struct}.

Разрешаването на имената до връзки се извършва \term{след} обхождането, когато всички
класове вече са известни, защото типът на член може да сочи към клас, деклариран по-късно
в същата единица за превод. Съответствието на име се прави първо по напълно квалифицирано
име, а ако то не съвпадне, по последната съставка. Двусмислено име не се разрешава до
предположение, а остава неразрешено и не поражда връзка.

\subsection{Заместващи заглавни файлове при извличането}

Извличането се изпълнява с \code{-nostdinc++} и с папката \code{stdstub/} като системна
папка за включване. Тя съдържа заглавни файлове само с декларации, заместващи тези от
стандартната библиотека, които заглавните файлове на проекта включват. Причината е
измерена, а не стилистична: една единица за превод, която включва действителните
\code{<memory>}, \code{<vector>} и \code{<string>}, извежда около 280 MB JSON, почти
изцяло вътрешности на стандартната библиотека, които никоя диаграма на класове не би
показала. Със заместващите файлове същата единица извежда под мегабайт. Те никога не се
компилират в инструмента и никога не се свързват, а извличането се изпълнява с
\code{-fsyntax-only}.

\subsection{Излъчватели}

Излъчвателите в двете нотации~\cite{plantuml_guide,mermaid_docs} споделят един обход на
модела и се различават само по низовете, които произвеждат за всеки елемент. Общата част
е изнесена в базов клас с виртуални методи за отделните елементи, а конкретните
излъчватели предефинират само тях. Редът на елементите и решението кое си заслужава да
бъде излъчено живеят в базовия клас, веднъж.

Това е и мястото, където се вижда цената на разликите между нотациите. Mermaid чете
ъгловите скоби като свой собствен ограничител за параметризация и използва тилда, така че
всяко изписване на тип се превежда на излизане. Нито една от двете нотации няма място за
етикетирана стойност или за списък на шаблонните параметри, затова те пътуват в
\term{бележка}. Бележката е механизмът на самия UML за пояснение върху елемент, тоест
това не е скрит канал, а предвиденият. Един признак остава без превод: Mermaid отбелязва
операция като абстрактна или като статична и няма нищо за операция, която е виртуална, но
не е чиста. Загубата е измерена, а не премълчана, и е отчетена в Раздел~\ref{sec:results}.

Четецът на диаграми приема диалекта, който самият Blueprint произвежда. Ограничението е
съзнателно и е заявено като такова: целта на четенето е да затвори цикъла и да даде на
проверяващия модул модел на документираното решение, а не да приеме всяка диаграма, която
човек би могъл да нарисува.

\subsection{Пораждаща посока}

Пораждащата посока е реализирана като две отделни преобразувания над един и същ модел.
Първото произвежда описание на релационна схема, второто заготовки на класове. И трите
стратегии за наследяване са реализирани, като изборът се прави от етикетирана стойност
върху базовия клас и само при липсата ѝ от подадената от командния ред стойност. При
стратегията с една таблица всяка колона, дошла от наследник, се записва като допускаща
празна стойност, защото ред от друг наследник я оставя празна, и се добавя колона
разграничител. При стратегията с таблица на клас първичният ключ на наследника е и външен
ключ към базата. При стратегията с таблица на конкретен клас абстрактният клас не получава
таблица, а наследените колони се копират надолу.

Породеният DDL носи името на класа в коментар пред всяка таблица и името на атрибута в
коментар след всяка колона, чието име се различава от неговото. Това не е украса.
Възстановяването на име на клас от име на таблица би означавало отгатване на английско
единствено число, а отгатването е точно това, което един пълен цикъл не бива да съдържа.

Заготовките се пораждат само за декларации, тоест инструментът не пише тела на операции и
не докосва съществуващ файл, който вече съдържа реализация. Породеният заглавен файл носи
и макросите за анотиране, така че повторното извличане от него дава съпоставим модел, с
което цикълът се затваря. Че породеният файл е действителен C++, се проверява в
демонстрацията чрез \code{clang -fsyntax-only} върху него, а не се твърди.

\subsection{Проверяващ модул}

Проверяващият модул сравнява модел, извлечен от кода, с модел, прочетен от диаграма, и
извежда разликите като списък. Видовете разлика са девет: добавен, премахнат и променен,
поотделно за класификатор, член и връзка. Съответствието на връзките се прави по двата
края и по етикета, така че промяна на вида или на кратността се отчита като промяна, а не
като премахване плюс добавяне. Разликата е съществена: два реда вместо един за една и съща
промяна биха направили отчета по-шумен точно там, където той трябва да е най-полезен.

Посоката на думите е фиксирана: първият подаден модел е кодът, вторият е решението, така
че клас, който е в кода, но не е в диаграмата, е \term{добавен}, а обратното е
\term{премахнат}. Процесът завършва с код 0 при съвпадение и с код 1 при разлика, което е
и условието проверката да може да прекъсне построяване, вместо да чака някой да я
прочете.

\subsection{Модулни тестове}

Модулните тестове са 58 в 7 групи и не използват външна рамка. Изпълнителят е един
заглавен файл с макроси за твърдение и регистър, попълван при статично инициализиране;
групите се регистрират в CTest поотделно с \code{add\_test}, така че CTest съобщава по
един ред на група, а всяка група изброява своите случаи. Тестовете на четеца работят
върху записан образец на синтактично дърво, в който присъстват само полетата, които
четецът действително чете, тоест образецът е едновременно вход на теста и заявление какво
четецът изисква от фронтенда.

\newpage

% ============================================================================
% IV. Методика и експериментални резултати.
% Двата файла са ЕДИН раздел: 05 описва плана, 06 дава числата. Само 05 отваря
% \section, а 06 продължава без нов заглавен ред и без \newpage между тях.
% ============================================================================
% III. МЕТОДИКА И ЕКСПЕРИМЕНТАЛНИ РЕЗУЛТАТИ, част 1: ПЛАНЪТ.
% ВНИМАНИЕ: тук няма и не трябва да има нито едно измерено число. Числата се
% появяват едва в chapters/06_results.tex, което е продължение на СЪЩИЯ раздел.
\section{Методика и експериментални резултати}
\label{sec:method}

\subsection{Методика}

Проверката отговаря на четири отделни въпроса: дали извлеченият модел отговаря на изходния текст, колко
струва извличането, дали пълният цикъл запазва модела, и дали проверяващият модул открива нарочно
внесено разминаване. Смесването им в едно число би скрило точно това, което трябва да бъде видяно.
Измерванията са на Windows 10.0.26200, AMD Ryzen 5 3600 с 6 ядра, построяване с GNU 15.2.0
(MinGW-w64) в режим Release, CMake 4.3.2 с Ninja 1.13.2, стандарт C++20, а фронтендът при извличането
е clang 22.1.0. Средата се извежда от измервателния скрипт, тъй като резултатите зависят от версията
на фронтенда. Всяко измерване следва чисто построяване, времената се повтарят по три пъти и се
съобщава \term{най-малкото}, а не средното: по-бавните изпълнения се различават по планиране на
процесите, което не е свойство на инструмента. Всичко се възпроизвежда с \code{cmake --build build
--target measure}, а четвъртият експеримент с целта \code{divergence}.

При \textbf{първия експеримент} очакваният модел за всяка конструкция е записан на ръка \term{преди}
инструментът да бъде пуснат върху нея, тоест сравнението не е между два изхода на самия инструмент;
отделно върху породените заготовки се пуска фронтендът с \code{-fsyntax-only}, защото пораждащата
посока пише код, който никой човек не е чел. \textbf{Вторият} използва породени единици с нарастващ
брой класове, всеки с по един представител на всяка форма на член от Табл.~\ref{tab:ownership}, и двата
действителни входа; фронтендът, разборът и обхождането се измерват поотделно, защото се променят
различно. \textbf{Третият} превръща текст в модел, модела в диаграма и диаграмата обратно в модел;
показателят е дял на \term{запазените елементи}, а не единица минус дял на разликите, защото когато цял
класификатор не се върне, членовете му изобщо не се сравняват. \textbf{Четвъртият} променя само едната
страна на съгласувана двойка, по една промяна наведнъж, като преди всяка промяна се проверява, че
недокоснатата двойка съвпада; показателите са два, откриваемост и брой съобщения извън очакваното,
защото инструмент, който съобщава за всичко, постига идеална откриваемост и е безполезен.

% IV. МЕТОДИКА И ЕКСПЕРИМЕНТАЛНИ РЕЗУЛТАТИ, част 2: ЧИСЛАТА.
% ВНИМАНИЕ: всяко число тук е измерено в средата, описана в 05_method.tex, и се
% възпроизвежда с посочената там команда. Клетка без измерване остава \TODO.
% NB: тук НЯМА \section. Файлът продължава раздела, открит в 05_method.tex.
\label{sec:results}

\subsection{Резултати}

\textbf{Вярност на извличането.}

Изпълнението е \code{ctest --test-dir build --output-on-failure}: 8 групи, 63 случая, всички
преминават, когато \code{clang++} е наличен; групата \code{extract} извежда \code{SKIP} и CTest я
отчита като пропусната, когато фронтендът липсва, а не като успешно преминала. Размерът на дървото и
на записания модел по-долу съдържат абсолютните пътища на
входните файлове, затова тези две колони се отместват с постоянна стойност при друг път; никоя
друга стойност не зависи от пътя.

\begin{table}[H]
\caption{Вярност на извличането по езикови конструкции}
\label{tab:extraction}
\centering
\setstretch{1.0}
\fontsize{10}{12}\selectfont
\begin{tabular}{@{}L{5.1cm}cL{5.1cm}c@{}}
\toprule
\textbf{Конструкция} & & \textbf{Конструкция} & \\
\midrule
Наследяване с достъп на базата & \cmark & Асоциативен контейнер & \cmark \\
Виртуално наследяване & \cmark & Статичен член, обхват на класа & \cmark \\
Множествено наследяване & \cmark & Чиста виртуална операция & \cmark \\
Член по стойност, композиция 1 & \cmark & Извеждане на интерфейс & \cmark \\
Притежаващ указател, композиция & \cmark & Шаблонен клас & \cmark \\
Споделен указател, агрегация & \cmark & Влагане в пространство от имена & \cmark \\
Указател без притежание & \cmark & Отсяване на подразбиращи се членове & \cmark \\
Член отправка, асоциация 1 & \cmark & Анотация за съхранение & \cmark \\
Контейнер от стойности, кратност * & \cmark & Отсяване по файл на произход & \cmark \\
Контейнер от указатели & \cmark & Породените заготовки са валиден C++ & \cmark \\
\bottomrule
\end{tabular}
\end{table}

Таблицата казва, че всяка изброена конструкция се извлича така, както е записано предварително. Тя
не казва, че покритието на езика е пълно: извън проверката остават приятелските декларации,
вложените шаблони на шаблони, наследяването от специализация на шаблон и всичко, което
препроцесорът е скрил при конкретния набор от ключове. Тези случаи не са проверени и затова не се
твърди нищо за тях.

\textbf{Цена на извличането.} Породените входове съдържат от 10 до 800 класа, всеки с по един представител на всяка форма на
член; съобщава се най-малкото от три повторения \tabref{tab:cost}.

\begin{table}[H]
\caption{Цена на извличането спрямо размера на единицата за превод}
\label{tab:cost}
\centering
\setstretch{1.0}
\fontsize{10}{12}\selectfont
\begin{tabular}{@{}rrrrrrrr@{}}
\toprule
\textbf{Класове} & \textbf{Изх. текст} & \textbf{Дърво} & \textbf{Възли} & \textbf{Елементи} &
\textbf{Фронтенд} & \textbf{Разбор} & \textbf{Обход} \\
 & \textbf{[B]} & \textbf{[B]} & & & \textbf{[s]} & \textbf{[s]} & \textbf{[s]} \\
\midrule
 10 &   3\,339 &      554\,127 &    3\,956 &    158 & 0.0631 & 0.0027 & 0.0004 \\
 25 &   8\,559 &      798\,297 &    5\,861 &    413 & 0.0688 & 0.0042 & 0.0010 \\
 50 &  17\,259 &   1\,206\,503 &    9\,036 &    838 & 0.0744 & 0.0067 & 0.0021 \\
100 &  34\,660 &   2\,024\,872 &   15\,386 & 1\,688 & 0.0932 & 0.0116 & 0.0055 \\
200 &  70\,360 &   3\,662\,672 &   28\,086 & 3\,388 & 0.1161 & 0.0215 & 0.0124 \\
400 & 141\,760 &   6\,945\,391 &   53\,486 & 6\,788 & 0.1790 & 0.0413 & 0.0323 \\
800 & 284\,560 &  13\,524\,380 &  104\,286 & 13\,588 & 0.2871 & 0.0815 & 0.0925 \\
\bottomrule
\end{tabular}
\end{table}

\begin{table}[H]
\caption{Двата действителни входа}
\label{tab:realinputs}
\centering
\setstretch{1.0}
\fontsize{10}{12}\selectfont
\begin{tabular}{@{}L{3.4cm}rrrrrr@{}}
\toprule
\textbf{Вход} & \textbf{Изх. текст} & \textbf{Дърво} & \textbf{Класи-} & \textbf{Връзки} &
\textbf{Модел} & \textbf{Общо} \\
 & \textbf{[B]} & \textbf{[B]} & \textbf{фикатори} & & \textbf{[B]} & \textbf{[s]} \\
\midrule
Анотиран пример & 3\,069 & 551\,549 & 6 & 8 & 6\,549 & 0.066 \\
Заглавни файлове на инструмента & 612 & 2\,952\,078 & 35 & 30 & 33\,366 & 0.108 \\
\bottomrule
\end{tabular}
\end{table}

От двете таблици следват четири неща. Първо, дървото расте линейно спрямо изходния текст, но с
коефициент около 48. Това е причината заместващите заглавни файлове да не са удобство, а условие: при
действителната стандартна библиотека същият коефициент дава стотици мегабайти. Второ, изпълнението на
фронтенда е най-скъпото до около 400 класа, но не расте пропорционално: при осемдесеткратно увеличение
на текста нараства 4.5 пъти, защото значителна част е постоянна цена за създаване на процес, докато
разборът на JSON расте практически линейно. Трето, обхождането до модел расте \term{по-бързо от
линейното}: между 400 и 800 класа входът се удвоява, а времето се увеличава 2.9 пъти, защото
разрешаването на име търси последователно в списък от всички класификатори. Поправката е очевидна,
речник вместо списък, и е записана като по-нататъшна работа, а не приложена, защото при действителните
входове разликата не е измерима. Четвърто, изходният текст на втория действителен вход е 612 байта,
защото файлът съдържа само включвания на деветте публични заглавни файла, а записаният модел е около 33
KB срещу близо 3 MB дърво. Това е количествената форма на твърдението, че моделът е абстракция: почти
цялото дърво е синтаксис, който диаграмата не показва.

\textbf{Запазване при пълен цикъл.}
\begin{table}[H]
\caption{Запазване на модела при пълен цикъл}
\label{tab:roundtrip}
\centering
\setstretch{1.0}
\fontsize{10}{12}\selectfont
\begin{tabular}{@{}L{3.7cm}L{2.2cm}rrrr@{}}
\toprule
\textbf{Вход} & \textbf{Път} & \textbf{Елементи} & \textbf{Запазени} & \textbf{Разлики} & \textbf{Дял} \\
\midrule
Анотиран пример & PlantUML & 41 & 41 & 0 & 1.000 \\
Анотиран пример & Mermaid  & 41 & 38 & 3 & 0.927 \\
Анотиран пример & схема    & 41 & 12 & 27 & 0.293 \\
Заглавни файлове & PlantUML & 226 & 226 & 0 & 1.000 \\
Заглавни файлове & Mermaid  & 226 & 225 & 1 & 0.996 \\
Заглавни файлове & схема    & 226 & 0 & 65 & 0.000 \\
\bottomrule
\end{tabular}
\end{table}

Пътят през PlantUML е без загуба и за двата входа. Това не беше вярно от самото начало: при първото
измерване се губеха имената на безименните параметри, защото излъчвателят ги записваше като
\code{: тип} и четецът прочиташе двоеточието обратно като част от типа. Дефектът беше открит именно от
този експеримент, а не от преглед на кода, и е поправен. Пътят през Mermaid губи по един елемент на
всяка операция, която е виртуална, но не е чиста: нотацията отбелязва операция като абстрактна или като
статична и няма трети знак. Загубата е свойство на нотацията, не на модела, и не подлежи на поправка.
Такива операции са три при анотирания пример и една при заглавните файлове, откъдето и разликата между
0.927 и 0.996, тоест делът зависи от това какъв дял от модела са виртуалните операции, а не от размера
му.

Пътят през схемата губи най-много, и това е очакваното: таблицата няма операции, няма видимост, няма
шаблони и няма класове, които не са обозначени за съхранение. При заглавните файлове се запазва нула,
защото нито един клас на инструмента не е обозначен. Нулата не е дефект, а вярно измерване на въпроса
колко от този модел изобщо може да бъде изразено като таблици. Общият извод е, че цикълът е затворен
през нотациите и е \term{частичен} през схемата: достатъчен, за да се отговори дали кодът и диаграмата
се разминават, и недостатъчен, за да се възстанови кодът от схемата. Второто и не се твърди.

\textbf{Откриваемост на внесено разминаване.}
\begin{table}[H]
\caption{Откриваемост на нарочно внесено разминаване}
\label{tab:divergence}
\centering
\setstretch{1.0}
\fontsize{10}{12}\selectfont
\begin{tabular}{@{}L{2.9cm}L{2.9cm}L{1.5cm}L{2.1cm}L{2.4cm}@{}}
\toprule
\textbf{Вид на промяната} & \textbf{Очакван вид съобщение} & \textbf{Открита} & \textbf{Общо съобщения} & \textbf{Извън очакваното} \\
\midrule
Премахнат атрибут & добавен член      & \cmark & 1 & 0 \\
Добавена операция & премахнат член    & \cmark & 1 & 0 \\
Сменена видимост  & променен член     & \cmark & 1 & 0 \\
Премахнат базов клас & добавена връзка & \cmark & 1 & 0 \\
Сменена кратност  & променена връзка  & \cmark & 2 & 1 \\
\bottomrule
\end{tabular}
\end{table}

И петте внесени разминавания са открити, при нула съобщения, които не отговарят на действително
разминаване. Второто число е важното: инструмент, който съобщава за всичко, би постигнал същата
откриваемост и би бил безполезен. Четири от петте случая дават точно по едно съобщение. Петият дава
две: при смяна на кратността от \code{0..1} на \code{1} се съобщава и за променена връзка, и за
променен атрибут. Второто не е грешно, защото кратността се пази на две места, върху връзката и върху
атрибута, който я е породил, и промяната на диаграмата променя и двете, тоест едно редактиране води
до две верни съобщения. Наблюдението обаче сочи истинско проектно решение: кратността е
\term{дублирана} в модела. Дублирането е удобно за излъчвателите и неудобно за проверяващия модул,
който брои една промяна два пъти; отстраняването му е въпрос на избор кой носител е първичен, и е
записано като по-нататъшна работа.

\newpage

% ============================================================================
% ЗАКЛЮЧЕНИЕ, без номер
% ============================================================================
% ЗАКЛЮЧЕНИЕ, без номер.
% По А.3: обобщение на характеристиките и особеностите на решението, предимствата и
% недостатъците му, резултатите от експериментите и тяхната инженерна трактовка, и
% предложения за по-нататъшна работа.
\section*{ЗАКЛЮЧЕНИЕ}
\label{sec:conclusion}
\addcontentsline{toc}{section}{ЗАКЛЮЧЕНИЕ}

В курсовия проект е проектиран и реализиран инструмент за двупосочно инженерство между
изходен код на C++, UML диаграма на класове и релационна схема на данни. Приетият подход
поставя в средата един вътрешен модел, независим както от езика на входа, така и от
формата на изхода, и свежда всяко преобразувание до преход между този модел и една външна
форма. Следствието е, че добавянето на нова нотация не засяга извличането, а промяна в
извличането се отразява едновременно във всички изходни формати.

Предимствата на решението произтичат от избора на компилаторен фронтенд за прочитане на
кода. Понятия като разрешено име, тип след заместване на шаблонните аргументи и
собственост върху обект са достъпни точно защото анализът е семантичен, а не текстов. Това
е и условието разграничението между агрегация и композиция да бъде правено по устройство,
а не по предположение. Второто предимство е в начина на достъп до фронтенда: той се
изпълнява като външен процес, а изводът му се разбира като текст, така че проектът няма
нито една задължителна външна зависимост и се построява само с компилатор за C++20 и
CMake. Ограниченията са в същата посока: инструментът вижда само това, което изводът на
фронтенда съдържа, а при това извличане стандартната библиотека е заместена с декларации,
тоест типове извън заявената папка остават непрозрачни атрибути, а не връзки.

Измереното от четирите експеримента дава следната картина. Всички 20 проверени езикови
конструкции се извличат така, както очакваният модел е записан предварително, а породените
от инструмента заготовки на класове са приети от компилатора. Цената на извличането е под
0.11 s за действителните входове, като изведеното синтактично дърво е около 48 пъти
по-голямо от изходния текст, а полученият от него модел е около 88 пъти по-малък от
дървото, тоест абстракцията е количествено измерима, а не само декларирана. Пълният цикъл
през PlantUML е без загуба и за двата входа, през Mermaid запазва 0.996 от елементите при
заглавните файлове на инструмента и 0.927 при анотирания пример, а през релационната схема
запазва 0.293 при анотирания пример и нула при заглавните файлове. И петте нарочно внесени
разминавания са открити, при нула съобщения, които не отговарят на действително
разминаване.

Три от тези числа заслужават трактовка, а не само отчитане. Първо, разликата между 1.000
и 0.996 не е дефект на модела, а свойство на нотацията: Mermaid отбелязва операция като
абстрактна или като статична и няма трети знак за операция, която е виртуална, но не е
чиста. Именно това е ползата от два излъчвателя вместо един. Второ, нулата през
релационната схема при заглавните файлове на инструмента е вярна: нито един клас на
инструмента не е обозначен за постоянно съхранение, така че въпросът „колко от този модел
може да бъде изразено като таблици“ има отговор нула. Трето, обхождането до модел расте
по-бързо от линейното, защото разрешаването на име търси последователно в списък от всички
класификатори; при 800 класа това вече е най-скъпата от трите съставки на извличането.

Наред с това самите експерименти откриха един действителен дефект, който преглед на кода
не беше открил: безименните параметри не преживяваха пълния цикъл, защото излъчвателят ги
записваше без име, а четецът прочиташе двоеточието обратно като част от типа. Дефектът е
поправен и е покрит с тест. Това е и най-прекият довод в полза на самия подход: щом и
двете посоки са автоматизирани, разминаването става измеримо, а не подозирано, и то важи
както за чужд код, така и за самия инструмент.

За по-нататъшна работа се очертават пет посоки. Първата е замяна на последователното
търсене при разрешаване на имената с речник, което премахва квадратичното поведение,
видимо над 400 класа. Втората е премахване на дублирането на кратността, която днес се пази
и върху връзката, и върху атрибута, поради което една промяна поражда две верни съобщения
вместо едно. Третата е разширяване на покритието на езика към конструкции, останали извън
проверката: приятелски декларации, наследяване от специализация на шаблон и вложени
шаблони. Четвъртата е свързване на проверката за разминаване с процеса на непрекъсната
интеграция, за което инструментът вече е готов, тъй като завършва с ненулев код при
разлика. Петата е допълване на модела с поведенческа информация, което е предпоставка за
извличане, а не само за излъчване, на диаграмите на случаите на употреба и на дейностите.

\newpage

% ============================================================================
% Списък на използваните литературни източници
% ============================================================================
\phantomsection
\addcontentsline{toc}{section}{ЛИТЕРАТУРА}
\bibliographystyle{ieeetr}
\bibliography{references}
\newpage

% ============================================================================
% Приложение
% ============================================================================
% ПРИЛОЖЕНИЯ, без номер на раздел.
% По А.3: изходен текст на програми, формати на структури от данни, екранни форми,
% получени резултати и др. Изходният текст се ВМЪКВА с \lstinputlisting от src/,
% а не се преписва тук: преписаният код се разминава с истинския още при първата
% поправка. Всички листинги от diagrams/ са ПОРОДЕНИ от инструмента и се
% възпроизвеждат с  cmake --build build --target demo
\section*{ПРИЛОЖЕНИЯ}
\label{sec:appendix}
\addcontentsline{toc}{section}{ПРИЛОЖЕНИЯ}

\subsection*{Приложение А. Формат на вътрешния модел}
\addcontentsline{toc}{subsection}{Приложение А. Формат на вътрешния модел}

Вътрешният модел е записваем във външен вид, за да може да бъде сравняван и съхраняван
между отделните стъпки на пълния цикъл. Записът е JSON с два масива на най-горно ниво,
\code{classifiers} и \code{relations}. Класификаторът носи неквалифицирано и квалифицирано
име, вид, признак за абстрактност, шаблонни параметри, стереотипи, етикетирани стойности,
атрибути, операции, изброители и мястото в изходния текст, откъдето е дошъл. Атрибутът
носи име, тип както е изписан в кода, видимост, обхват, кратност и собствени етикетирани
стойности. Операцията носи име, тип на резултата, параметри, видимост, обхват и признаци за
виртуалност, чистота и константност. Връзката носи двата края, вид, етикет, кратност,
признак за виртуална база и видимост на базата.

Класификаторите се свързват по име, а не по указател. Изборът е продиктуван от изискването
моделът да може да бъде записан и прочетен обратно, а също и да бъде сравняван с друг модел
по стойност: имената преживяват запис във файл, указателите не.

\lstinputlisting[language=C++, caption={Дефиниция на вътрешния модел, \code{include/blueprint/model.hpp}}, label={lst:model}]{../include/blueprint/model.hpp}

\subsection*{Приложение Б. Примерен вход и получените от него диаграми}
\addcontentsline{toc}{subsection}{Приложение Б. Примерен вход и получените диаграми}

Приложението съдържа малката единица за превод, върху която е проведен пълният цикъл. Тя
обхваща наследяване, абстрактен базов клас, асоциация, агрегация, композиция и кратност
много, като всяка от тези конструкции се среща в нея точно по веднъж, за да сочи един
евентуален отказ към едно правило, а не към сплетени правила.

\lstinputlisting[language=C++, caption={Анотиран примерен вход, \code{examples/library.hpp}}, label={lst:library}]{../examples/library.hpp}

Следват двата текста, получени от този вход. И двата са породени от един и същ модел с
\code{blueprint diagram} и се различават само по нотацията.

\lstinputlisting[language={}, caption={Породена диаграма на класове в PlantUML}, label={lst:puml}]{../diagrams/library.class.puml}

\lstinputlisting[language={}, caption={Породена диаграма на класове в Mermaid}, label={lst:mmd}]{../diagrams/library.class.mmd}

Поведенческите модели се превеждат по същия начин. По-долу е диаграмата на случаите на
употреба в PlantUML и диаграмата на дейностите в Mermaid, тоест същите два модела, които
Фиг.~\ref{fig:usecase} и Фиг.~\ref{fig:activity} показват.

\lstinputlisting[language={}, caption={Породена диаграма на случаите на употреба в PlantUML}, label={lst:ucpuml}]{../diagrams/blueprint.usecase.puml}

\lstinputlisting[language={}, caption={Породена диаграма на дейностите в Mermaid}, label={lst:actmmd}]{../diagrams/roundtrip.activity.mmd}

\subsection*{Приложение В. Породена релационна схема и обратно породени заготовки}
\addcontentsline{toc}{subsection}{Приложение В. Породена релационна схема}

Схемата по-долу е получена от модела на Приложение Б при стратегия \term{една таблица за
йерархията}. Класовете \code{Book} и \code{Journal} са свити в таблицата \code{items},
която получава колона разграничител, а колоните им допускат празна стойност. Асоциацията с
кратност много между \code{Catalogue} и \code{LibraryItem} поражда свързваща таблица.
Името на класа пътува в коментар пред всяка таблица, а името на атрибута в коментар след
колоната, чието име се различава от неговото, за да не се налага при обратния преход да се
отгатва единствено число на английска дума.

\lstinputlisting[language=SQL, caption={Породена схема при стратегия „една таблица за йерархията“}, label={lst:sql}]{../diagrams/library.single-table.sql}

Обратният преход от схема към код затваря цикъла. Породеният заглавен файл съдържа само
декларации и носи макросите за анотиране, така че повторното извличане от него дава
съпоставим модел. Че той е действителен C++, се проверява в демонстрацията с
\code{clang -fsyntax-only}, а не се твърди.

\lstinputlisting[language=C++, caption={Заготовки на класове, породени обратно от схемата}, label={lst:skeleton}]{../diagrams/library_generated.hpp}

\subsection*{Приложение Г. Ключови места от изходния текст}
\addcontentsline{toc}{subsection}{Приложение Г. Ключови места от изходния текст}

Пълният изходен текст е в хранилището. Тук се вмъкват двете места, на които се взема
същинското решение, а не се пренася информация.

Първото е правилото, което превръща изписването на един член във вид на връзка. То е
съзнателно събрано в една функция, за да може да бъде прочетено наведнъж и оспорено
наведнъж.

\lstinputlisting[language=C++, firstline=133, lastline=212, firstnumber=133, caption={Правилото за собствеността, \code{src/reader.cpp}}, label={lst:ownership}]{../src/reader.cpp}

Второто е броенето на запазените елементи, което е числителят на дела от
Табл.~\ref{tab:roundtrip}. То се брои поелементно, а не се извежда от броя на разликите,
и коментарът в кода казва защо.

\lstinputlisting[language=C++, firstline=116, lastline=142, firstnumber=116, caption={Броене на запазените елементи, \code{src/diff.cpp}}, label={lst:preserved}]{../src/diff.cpp}


\end{document}
